% Einheit 2 von 5 – Prozentsatz berechnen (bank/prozentrechnung/e2.jsonl, zusammenbau v0.3)
%% TODO Zweigzeile Teil 1: was hier gelernt wird (Satz in Schülersprache aus dem Einheitstitel) – Einheit 2
\einheitenkopf[e2]{Einheit 2 von 5 · Prozentsatz berechnen}
\zweigzeile{ab Kl. 7 · P10 oft}
\setcounter{aufgabe}{11}

%% TODO Titel: Ich-kann-Satz fehlt in der Bank (Platzhalter: Kettenname) – Nr. 12
%% TODO Anweisung: ein Satz über der ersten Teilaufgabe fehlt in der Bank – Nr. 12
\begin{aufgabe}{Was ist das Ganze?}
\begin{teile}
\teil In der Klasse sind $28$ Kinder, $7$ davon haben einen Hund. – Was ist das Ganze?
\end{teile}
\end{aufgabe}

%% TODO Titel: Ich-kann-Satz fehlt in der Bank (Platzhalter: Kettenname) – Nr. 13
%% TODO Anweisung: ein Satz über der ersten Teilaufgabe fehlt in der Bank – Nr. 13
\begin{aufgabe}{Streifen einteilen}
\begin{teile}
\teil In $10$-\%-Schritte einteilen.

\streifenleer[0]
\end{teile}
\end{aufgabe}

%% TODO Titel: Ich-kann-Satz fehlt in der Bank (Platzhalter: Kettenname) – Nr. 14
%% TODO Anweisung: ein Satz über der ersten Teilaufgabe fehlt in der Bank – Nr. 14
\begin{aufgabe}{Prozentsatz}
%% TODO \rechenplatz{n}: Schrittzahl des Grundfalls fehlt in der Bank (nur bei fester Schreibform, 2.3 b)
\begin{teile}
\teil Streifen: $50$ Kinder, grau: $10$ Kinder – wie viel Prozent?

\streifenfeld{20}{$0$}{$50$ Kinder}
\teil $17$ von $100$ Schülern spielen ein Instrument – wie viel Prozent? \leerfeld[\%]
\teil $19$ von $50$ Kindern – wie viel Prozent? \leerfeld[\%]
\teil $27$ von $60$ Minuten Training sind Laufen – wie viel Prozent? \leerfeld[\%]
\teil $17$ von $24$ Kindern waren im Schwimmbad – wie viel Prozent? Runde auf eine Stelle. \leerfeld[\%]
\teil Im Chor singen $14$ Mädchen und $6$ Jungen – wie viel Prozent sind Jungen? \leerfeld[\%]
\teil Jacke: vorher $80\,€$, jetzt $60\,€$ – wie viel Prozent Rabatt? \leerfeld[\%]
\teil Bei einer Tombola gibt es $21$ Nieten und $3$ Hauptgewinne. Wie viel Prozent der Lose sind Hauptgewinne? \hfill (P10 2018 OS) \\ \leerfeld[\%]
\teil Eine Stadt hat $40\,000$ Einwohner, $1\,080$ sind in diesem Jahr zugezogen. Wie viel Prozent sind das? \hfill (P10 2014 OS) \\ \leerfeld[\%]
\teil Körpergrößen der $13$ Spieler einer Mannschaft in cm: $172$, $185$, $169$, $190$, $178$, $181$, $166$, $188$, $175$, $183$, $170$, $192$, $177$. Wie viel Prozent sind größer als $180$ cm? Runde auf eine Stelle. \hfill (P10 2025 OS)
\end{teile}
\end{aufgabe}

%% TODO Titel: Ich-kann-Satz fehlt in der Bank (Platzhalter: Kettenname) – Nr. 15
%% TODO Anweisung: ein Satz über der ersten Teilaufgabe fehlt in der Bank – Nr. 15
\begin{aufgabe}{Umkehrung: zu einem Prozentsatz ein Zahlenpaar angeben}
\begin{teile}
\teil $30\,\%$ – Teil und Ganzes? \leerfeld von \leerfeld
\end{teile}
\end{aufgabe}

%% TODO Titel: Ich-kann-Satz fehlt in der Bank (Platzhalter: Kettenname) – Nr. 16
%% TODO Anweisung: ein Satz über der ersten Teilaufgabe fehlt in der Bank – Nr. 16
\begin{aufgabe}{Prozentsatz – Fehler finden, Begründen, Darstellungswechsel, Anwendung}
\begin{teile}
\teil Kim, „$12$ von $48$ Kindern“: \rechnung{48 : 12 &= 4 = 400\,\%} Fehler? Wie heißt es richtig?
\teil Warum teilt man beim Prozentsatz durch das Ganze und nicht durch den Teil?
\teil Ganzes $40$ kg, Teil $10$ kg – im Streifen einzeichnen, Prozentsatz ablesen.

\streifenleer[4]
\leerfeld[\%]
\teil Handyvertrag mit $20$ GB im Monat. Am 20. Tag sind $17$ GB verbraucht. Wie viel Prozent sind noch übrig? Reicht der Rest bei gleichem Verbrauch für die letzten $10$ Tage?
\end{teile}
\end{aufgabe}
